\subsection{Etap Code}
Faza wytwarzania (\textit{Code}) obejmuje implementację wymagań przy zachowaniu standardów jakości kodu i spójności historii zmian. Fundamentem tego etapu jest rozproszony system kontroli wersji Git, a konwencje nazewnictwa zatwierdzeń są wymuszane na poziomie potoku CI.

\subsubsection{System kontroli wersji: GitHub}
Jako system kontroli wersji i platformę współpracy wybrano \textbf{GitHub}, oparty na rozproszonym systemie Git. Jest on centralnym punktem prawdy (\textit{Single Source of Truth}) dla całego monorepozytorium.

\textbf{Standard Conventional Commits:}
Zgodnie z decyzją projektową \textit{ADR-0004}, wprowadzono standard \textbf{Conventional Commits}. Polega on na nadawaniu wiadomościom zatwierdzeń ustrukturyzowanej formy (np. \texttt{feat:}, \texttt{fix:}, \texttt{chore:}).
\begin{itemize}
    \item \textbf{Uzasadnienie:} Standaryzacja pozwala na automatyczne generowanie list zmian (\textit{changelogs}) oraz automatyzację wersjonowania semantycznego (\textit{Semantic Versioning}).
    \item \textbf{Rola w CI/CD:} Czytelna historia zmian umożliwia narzędziom automatyzującym precyzyjne określenie, które komponenty systemu wymagają przebudowania i wdrożenia po każdej zmianie w monorepozytorium.
\end{itemize}

Poniższy przykład ilustruje kompletną wiadomość zatwierdzenia zgodną ze standardem. Nagłówek zawiera typ, zakres oraz zwięzły opis zmiany, a stopka (\textit{footer}) powiązuje commit z konkretnym zadaniem w systemie Jira:

\begin{lstlisting}[caption={Przykład wiadomości zatwierdzenia zgodnej ze standardem Conventional Commits}, label={lst:conventional-commit}]
feat(helm): create umbrella chart for all 7 services

Aggregates all microservice charts under a single parent
chart, enabling one-command deployment of the full stack.

Resolves: INZ-9
\end{lstlisting}

\subsubsection{Wytwarzanie oprogramowania: Stos technologiczny}
Proces kodowania realizowany jest przy użyciu nowoczesnych środowisk programistycznych (\textit{IntelliJ IDEA}, \textit{VS Code}), a sama implementacja opiera się na technologiach wspierających paradygmat mikroserwisów.

\textbf{Backend (Java i Quarkus -- tryb GraalVM Native Image):}
Mikroserwisy backendowe zbudowano z użyciem frameworka \textit{Quarkus} skompilowanego do natywnych plików wykonywalnych za pomocą \textit{GraalVM Native Image} (kompilacja \textit{Ahead-of-Time}, AOT). W trybie natywnym Quarkus analizuje ścieżki wykonania kodu w czasie kompilacji i wyklucza nieużywane klasy, co skutkuje:
\begin{itemize}
    \item \textbf{Błyskawicznym startem} -- aplikacja jest gotowa do obsługi żądań w czasie poniżej 100 ms, co skraca czas wznowienia po awarii i przyspiesza skalowanie poziome w Kubernetes.
    \item \textbf{Niskim zużyciem pamięci} -- brak środowiska uruchomieniowego JVM zmniejsza footprint kontenera, co bezpośrednio przekłada się na efektywność kosztową w środowiskach chmurowych.
    \item \textbf{Minimalnym obrazem kontenera} -- obraz oparty na \texttt{ubi-minimal} zawiera wyłącznie system operacyjny i skompilowany binarny, bez JDK ani bibliotek środowiska uruchomieniowego.
\end{itemize}
Wstrzykiwanie zależności realizowane jest przez mechanizm CDI Quarkusa (\textit{ArC}), który w trybie AOT przetwarza wszystkie adnotacje w czasie kompilacji, a nie w czasie uruchomienia.

\textbf{Frontend (Vue 3 + Vite):}
Zgodnie z decyzją \textit{ADR-0014}, warstwę prezentacji zbudowano z użyciem frameworka \textit{Vue 3} i narzędzia budującego \textit{Vite}. Kod pisany jest w \textit{TypeScript}, a poprawność typów weryfikuje \textit{vue-tsc} przed każdym buildem produkcyjnym. Komponenty interfejsu komunikują się z mikroserwisami backendowymi przez REST API.

\textbf{Pryncypia Twelve-Factor App:}
Wszystkie komponenty systemu są projektowane zgodnie z zasadami \textit{Twelve-Factor App}. Kluczowe aspekty to:
\begin{itemize}
    \item \textbf{Konfiguracja:} Cała konfiguracja specyficzna dla środowiska jest wstrzykiwana przez zmienne środowiskowe, a nie zaszyta w kodzie.
    \item \textbf{Procesy:} Aplikacje są bezstanowe (\textit{stateless}), co umożliwia ich łatwe skalowanie w poziomie wewnątrz klastra Kubernetes.
    \item \textbf{Logowanie:} Logi są traktowane jako strumień zdarzeń wysyłany na standardowe wyjście, co umożliwia ich centralizację przez narzędzia observability.
\end{itemize}
